\documentclass{article}
\usepackage{graphicx}
\usepackage[spanish,activeacute,es-tabla]{babel}
\usepackage[utf8]{inputenc} % mac utf8
\usepackage{amsfonts}
\usepackage{minted}
\usepackage{ucs}

\begin{document}

\title{\textit{Teoría de la Información y Codificación}\\ \textbf{Práctica 0: Toma de Contacto con la Asignatura}}
\author{José A. Montenegro Montes}

\maketitle




\section{Enunciado}


En esta tarea vamos a codificar cinco tipos de algoritmos de detección de errores: NIF, ISBN, ISBN13, Código de producto (UPC) y Código Control número de cuenta bancaria (opcional). 

Los algoritmos para crear los códigos de control mencionados son proporcionados en las transparencias de la asignatura y adjuntos en la última sección de las prácticas, menos el Código Control número de cuenta bancaria.

La práctica consta de la creación de una clase abstracta de codificación, y la implementación de los cuatros algoritmos solicitados (y uno opcional).

Los primeros tres métodos están implementados en el código aportado. El alumno debe revisar que la implementación cumple con la definición de los algoritmos e implementar el último método($corregirDatos$).


\begin{itemize}
\item Creación de un código control para un mensaje dado ($generarCodigoControl$). 
\item Verificación mensaje es correcto según el código control ($verificar$).
\item Corrección del código control erróneo de un mensaje dado ($corregirControl$).
\item Corrección del bit erróneo de un mensaje (un único error posible), dado un código control correcto ($corregirDatos$).
\end{itemize}

En el caso concreto del ISBN-10 sabemos que el código de control es mod 11, que ocurre si el código de control que obtenemos es 10. ¿Cómo lo incluimos en el código?.


\section{Conclusiones}

El objetivo de esta práctica es tomar contacto con este tipo de algoritmos y observar las ventajas y desventajas que nos presentan estos algoritmos. 

Se abre el debate si son códigos correctores o simplemente códigos detectores de errores. 

A lo largo del curso veremos una variedad de algoritmos tanto para la detección como para la corrección de errores, cuya calidad puede medirse en los número de bits que detecta / corrige.

\section{Código}


\subsection*{Clase Codificación } 

\inputminted[mathescape,
               ,
               numbersep=5pt,
               frame=lines,
               framesep=2mm]{java}{practica0/Codificacion.java}

               %%%%%%%%%%%%%%%%%%%%%%%
 
\subsection*{Clase ISBN10 Java.} 
\inputminted[mathescape,
               ,
               numbersep=5pt,
               frame=lines,
               framesep=2mm]{java}{practica0/ISBN.java}
               
%%%%%%
               %%%%%%%%%%%%%%%%%%%%%%%
\subsection*{Clase ISBN13 Java.} 

\inputminted[mathescape,
               ,
               numbersep=5pt,
               frame=lines,
               framesep=2mm]{java}{practica0/ISBN13.java}
               
%%%%%%


%%%%%%
               %%%%%%%%%%%%%%%%%%%%%%%
\subsection*{Clase NIF Java.} 


\inputminted[mathescape,
               ,
               numbersep=5pt,
               frame=lines,
               framesep=2mm]{java}{practica0/NIF.java}
               
%%%%%%

%%%%%%
               %%%%%%%%%%%%%%%%%%%%%%%
\subsection*{Clase UPC Java.} 

\inputminted[mathescape,
               ,
               numbersep=5pt,
               frame=lines,
               framesep=2mm]{java}{practica0/UPC.java}

%%%%%%

%%%%%%
               %%%%%%%%%%%%%%%%%%%%%%%
\subsection*{Clase Main Java.} 


\inputminted[mathescape,
               ,
               numbersep=5pt,
               frame=lines,
               framesep=2mm]{java}{practica0/Main.java}

               
%%%%%%


\end{document}
